The model draws on three bodies of work that rarely cite one another. We review what each establishes, note what it leaves open, and close with the connection none of them makes.

\subsection{How Newcomers Learn Ethical Norms}

Organizational socialization research treats entry as the period in which newcomers learn how things are done. \citet{vanmaanen1979toward} treated socialization as the way people come to acquire the social knowledge and skills that an organizational role demands, and they classified the tactics organizations use to structure it. Institutionalized tactics, in which newcomers move through shared, structured experiences, were later associated with lower role ambiguity and stronger identification~\citep{ashforth1996socialization}. A meta-analysis identified role clarity, self-efficacy, and social acceptance as indicators of newcomer adjustment that link socialization tactics and information seeking to later outcomes~\citep{bauer2007newcomer}.

One of the tactic dimensions in Van Maanen and Schein's typology is especially relevant here. They contrasted serial socialization, in which experienced members who have held a role groom newcomers to take it on, with disjunctive socialization, in which no such predecessors are available and newcomers have to work out the role largely on their own. Serial tactics give newcomers models to observe. The distinction was drawn with the availability of predecessors in mind, not physical distance, but it suggests a useful way to describe what distributed entry may do: newcomers who formally have experienced colleagues may in practice see so little of them that their situation comes to resemble disjunctive socialization.

Several strands of this research point to observation and interaction as the main channels. Newcomers make sense of the surprises of entry partly through the interpretations insiders offer~\citep{louis1980surprise}, and the frequency of interaction with insiders has been proposed as a main determinant of how quickly socialization proceeds~\citep{reichers1987interactionist}. In a study of how newcomers acquire different kinds of information, normative information was obtained mainly by observation and mainly from peers~\citep{morrison1993newcomer}.

Ethical norms are learned through the same channels. Early work found that subordinates' behavior came to resemble that of superiors they regarded as competent and successful~\citep{weiss1977subordinate}. Ethical leadership research built on this idea, defining ethical leaders in part as role models whose conduct, communication, and use of rewards and discipline show followers what is expected~\citep{brown2005ethical,brown2006ethical}. Evidence on ethical leadership is consistent with this account. A meta-analysis found that ethical leadership was positively associated with followers' ethical behavior~\citep{bedi2016meta}, and a multilevel study found that ethical leadership by top management was related to group outcomes through ethical leadership by supervisors~\citep{mayer2009how}. Peers matter too: employees' antisocial behavior has been found to track that of their coworkers~\citep{robinson1998monkey}.

Behavioral ethics research has examined influences on ethical conduct at several levels~\citep{trevino2006behavioral}, and meta-analytic evidence links features of the organizational environment, alongside characteristics of individuals and of the moral issue, to unethical choices~\citep{kishgephart2010bad}. Within that environment, researchers distinguish formal elements, such as codes, training, and official sanctions, from informal elements carried by everyday interaction~\citep{tenbrunsel2003building}, and they separate ethical culture from ethical climate as related but distinct features of the ethical context~\citep{trevino1998ethical}. In that study, a culture-based dimension of the ethical context was more strongly associated with observed unethical conduct in organizations that had codes of conduct, while climate-based dimensions mattered more in organizations without codes. Having a code, in other words, did not settle what was practiced. This distinction is central to our argument. Formal elements can be delivered to a newcomer anywhere. Informal elements, which tell a newcomer what is actually done, are the ones this literature ties most closely to interaction and observation.

What this body of work does not examine is what happens to ethical learning when the channels it relies on narrow. Its studies were conducted, and its theories developed, largely under the assumption that newcomers share space with the people they learn from.

\subsection{What Working at a Distance Changes}

Research on telework predates the recent spread of remote work and gives a mixed but informative picture. A meta-analysis found small, mostly favorable effects of telecommuting on outcomes such as perceived autonomy and work-family conflict, with the exception that high-intensity telecommuting was associated with worse coworker relationships~\citep{gajendran2007good}. A later review reached a similar overall conclusion while emphasizing that effects depend on how and how much telework is used~\citep{allen2015how}. Professional isolation was more harmful to teleworkers' performance the more time they spent teleworking, and less harmful the more face-to-face interaction they had~\citep{golden2008impact}.

The move to widespread working from home in 2020 raised these questions at a larger scale~\citep{kniffin2021covid}. In one large firm, the shift to firm-wide remote work made employees' collaboration networks more static and siloed and moved communication toward asynchronous media~\citep{yang2022effects}. A qualitative study of newcomers onboarded remotely found that, although they became proficient at their jobs, they had fewer interactions with insiders, experienced ambiguity about their membership, and found insider support insufficient~\citep{paik2026membership}. The same newcomers compensated with proactive behavior, seeking out the contact that did not come to them. That finding matters for our purposes because it implies that distant newcomers learn what they think to ask about, which is a narrower set of things than what co-located newcomers happen to see.

The common thread in these studies is that distance removes the unplanned and ambient part of work more than its scheduled part. Planned meetings continue, while the chance encounters, overheard conversations, and visible reactions that fill the time between them become rarer. Much of socialization research suggests that newcomers learn a great deal from exactly this unplanned part of work~\citep{reichers1987interactionist}.

Digital tools can restore some visibility. An enterprise social platform that made employees' messages visible to colleagues improved observers' knowledge of who knew what and who knew whom~\citep{leonardi2015ambient}. That finding concerns knowledge of expertise and connections, not ethical norms, but it shows that mediated channels can carry some of the ambient information that co-presence supplies.

In our reading, much of the earlier telework research also studied employees who had already been socialized in their organizations before they began working remotely. Newcomers who start at a distance are a different case, because they have no stock of previously observed norms to draw on. This literature has concentrated on performance, relationships, well-being, and adjustment. We found no study that examines how distance changes what newcomers learn about ethical norms.

\subsection{Generations, Cohorts, and Conditions of Entry}

The sociological concept of a generation goes back to~\citet{mannheim1952problem}, who argued that people born in the same period share a location in history but become a generation in a fuller sense only through shared formative experiences. Mannheim also expected people within the same broad cohort to respond to those experiences in different ways, so that a single cohort could contain several distinct groupings. His account therefore gives little support to the idea that everyone born within a span of years shares a common outlook. Workplace research has mostly used birth-year cohorts instead. Reviews of that research describe the evidence for generational differences in work values and attitudes as mixed, largely descriptive, and hard to separate from age and period effects~\citep{parry2011generational,lyons2014generational}, and a meta-analysis found differences in job satisfaction, commitment, and turnover intentions to be small~\citep{costanza2012generational}. One of these reviews concluded that the field had done much more to describe differences than to explain why a generation should differ at all~\citep{lyons2014generational}.

The field remains divided. One focal article argued that there is little solid evidence for generationally based differences at work and little theory explaining why they should exist~\citep{costanza2015generationally}, while a commentary defended generational differences as real and useful~\citep{campbell2015generational}. Critics have proposed replacing generational explanations with a lifespan developmental perspective~\citep{rudolph2017considering} and have challenged a series of common claims about generations~\citep{rudolph2021generations}.

A different line of work keeps the idea of generations but changes its basis. \citet{joshi2010unpacking} distinguished age-based generational identities from cohort-based ones formed when people enter an organization or experience at the same time. That distinction, together with Mannheim's emphasis on formative experience, opens a route that the value-difference debate has not taken: asking whether the conditions under which a cohort entered working life left a lasting mark on what it learned. The route does not require settling the debate between critics and defenders of generational research. It requires only that entry conditions can be measured and that their effects can be compared across newcomers of different ages, which is the approach taken in our final proposition.

\medskip

Taken together, these literatures leave a specific gap. Socialization and ethical leadership research explain how ethical norms are learned through observation but assume co-presence. Telework and remote-onboarding research documents the loss of interaction but not its ethical consequences. Imprinting research shows that conditions at entry leave durable traces, but it has studied performance and professional orientation~\citep{tilcsik2014imprint,azoulay2017social} rather than ethics. Generational research has focused on value levels. No existing model links distributed entry to the transmission and persistence of enacted ethical norms, and none considers that distance might weaken the transmission of unethical norms as well as ethical ones. The closest studies each supply one piece: remote-onboarding research the loss of contact with insiders, socialization research the importance of observation for normative learning, and imprinting research the durability of what is learned at entry. The model developed below joins these pieces around ethical norms.
